\documentclass[
  spanish,
  letterpaper, oneside
]{article}

\def\documenttitle {Bienvenida, introducción y objetivo de Garantías Constitucionales}
\def\documentsubtitle {Encuadre inicial, presentación docente y unidad didáctica}
\def\documentsubject {Licenciatura en Derecho}

\def\documentauthor {Martin Jonathan de la Cruz}
\def\coursename {Garantías Constitucionales}
\def\coursecode {025}

\def\universityname {Universidad Abierta y a Distancia de México}
\def\universityfaculty {Licenciatura en Derecho}
\def\universitydepartment {Garantías Constitucionales}
\def\universitydepartmentimage {departamentos/UnADM}
\def\universitydepartmentimagecfg {height=1.57cm}
\def\universitylocation {Roma Norte, Ciudad de México}

\def\authortable {
  \begin{tabular}{ll}
    \textbf{Alumno:}
    & \begin{tabular}[t]{l}
      Martin Jonathan de la Cruz \\
    \end{tabular} \\
    Matricula: & ES2611202040 \\
    Figura docente: & Aarón López Pérez \\
    Semestre/Bloque: & 2 / 1 \\
    & \\
    \multicolumn{2}{l}{Fecha de realizacion: \today} \\
    \multicolumn{2}{l}{\universitylocation}
  \end{tabular}
}

\input{template}
\usepackage{helvet}
\renewcommand{\familydefault}{\sfdefault}
\usepackage{array}
\usepackage{longtable}
\usepackage{pdflscape}
\usepackage{tikz}
\usetikzlibrary{arrows.meta,positioning,calc,matrix,fit,shapes.geometric,shadows.blur}
\setcitestyle{authoryear,open={(},close={)}}

\def\coverwatermarkenabled {true}
\def\coverwatermarkimage {img/departamentos/UnADM.pdf}
\def\coverwatermarkopacity {0.16}
\def\coverwatermarkwidth {0.72\paperwidth}
\def\coverwatermarkxshift {0cm}
\def\coverwatermarkyshift {-0.35cm}

\newcommand{\insertcoverwatermark}{%
  \ifthenelse{\equal{\coverwatermarkenabled}{true}}{%
    \AddToShipoutPictureBG*{%
      \begin{tikzpicture}[remember picture,overlay]
        \node[
          opacity=\coverwatermarkopacity,
          inner sep=0pt,
          xshift=\coverwatermarkxshift,
          yshift=\coverwatermarkyshift
        ] at (current page.center) {%
          \includegraphics[width=\coverwatermarkwidth]{\coverwatermarkimage}%
        };
      \end{tikzpicture}%
    }%
  }{}%
}

\begin{document}

\insertcoverwatermark
\templatePortrait
\templatePagecfg
\onehalfspacing

\begin{abstractd}
  Este documento integra las notas limpias de bienvenida, introducción y objetivo
  de la asignatura de \textit{Garantías Constitucionales}. Su propósito es
  conservar, en formato académico y verificable, el mensaje inicial dirigido al
  grupo, la presentación del docente, los conceptos clave, la relevancia de la
  materia para la abogacía y la ruta de análisis de derechos humanos, garantías
  de protección y medios de control constitucional.
\end{abstractd}

\templateIndex
\templateFinalcfg

\section{Bienvenida}

Estimadas y estimados estudiantes:

Me es muy grato comunicarles que iniciaremos la asignatura de \textit{Garantías Constitucionales}. Sean todas y todos bienvenidos. Quedaré muy agradecido de interactuar con ustedes y conocernos en la sesión sincrónica. Todos sus comentarios e inquietudes serán bienvenidos y escuchados en un marco de respeto.

\section{Presentación del docente}

El docente de la asignatura se presenta como \textbf{Aarón López Pérez}. Es abogado y cuenta con posgrados de Especialidad y Maestría realizados en la Facultad de Derecho de la Universidad Nacional Autónoma de México (UNAM). En su mensaje inicial expresa que será un placer coincidir con el grupo para aprender conjuntamente a partir de las opiniones de las y los estudiantes.

\section{Introducción a la unidad didáctica}

La asignatura de \textit{Garantías Constitucionales} es fundamental para la comprensión y el ejercicio de la abogacía. Abarca, por un lado, los derechos de las personas y, por otro, su defensa y protección, con fundamento en la Constitución mexicana. Por ello, el arranque de la materia debe incluir un diagnóstico de los conocimientos previos que permita realizar una sesión de retroalimentación o nivelación del pensamiento jurídico.

\section{Objetivo de la asignatura}

El objetivo de la unidad consiste en conocer y analizar los derechos humanos, las garantías para su protección y los medios de control constitucional previstos en la Constitución Política de los Estados Unidos Mexicanos, con el fin de abordar profesionalmente las problemáticas relacionadas con la defensa de los derechos humanos y promover la justicia y el desarrollo social.

\section{Conceptos clave de inicio}

\begin{longtable}{>{\bfseries}p{0.30\textwidth}p{0.62\textwidth}}
\toprule
Concepto & Sentido dentro de la asignatura \\
\midrule
Garantías Constitucionales & Asignatura orientada al estudio de los derechos de las personas y de los mecanismos para su defensa y protección con base en la Constitución mexicana. \\
Sesión sincrónica & Espacio previsto para interactuar y conocerse entre docente y estudiantes. \\
Marco de respeto & Ambiente en el que se recibirán y escucharán comentarios e inquietudes. \\
Posgrados & Estudios de Especialidad y Maestría con los que cuenta el docente. \\
Derechos humanos & Derechos cuyo análisis y defensa forman parte central del objetivo de la asignatura. \\
Medios de control constitucional & Mecanismos previstos en la Constitución Política de los Estados Unidos Mexicanos para la protección y defensa de derechos. \\
Diagnóstico de conocimientos previos & Evaluación inicial para organizar retroalimentación o nivelación del pensamiento jurídico. \\
Pensamiento jurídico & Base de comprensión que permite pasar de hechos a problemas jurídicos, normas aplicables, argumentos y vías de protección. \\
\bottomrule
\end{longtable}

\section{Ruta inicial de aprendizaje}

Para convertir la bienvenida en una experiencia formativa verificable, la materia puede iniciar con una ruta mínima de trabajo:

\begin{enumerate}
  \item Diagnosticar conocimientos previos mediante preguntas o microcasos no sancionadores.
  \item Distinguir entre derecho humano, garantía de protección y medio de control constitucional.
  \item Analizar casos breves a partir de hechos, derecho afectado, autoridad o acto, parámetro constitucional y remedio posible.
  \item Validar las soluciones con fundamento constitucional y criterios vigentes cuando corresponda.
  \item Mantener un marco de respeto que permita preguntar, disentir y argumentar con cuidado.
\end{enumerate}

\section{Términos relacionados}

\begin{itemize}
  \item Bienvenida.
  \item Presentación del docente.
  \item Estudiantes.
  \item Abogado.
  \item Facultad de Derecho.
  \item UNAM.
  \item Constitución mexicana.
  \item Defensa de los derechos humanos.
  \item Protección de derechos.
  \item Justicia y desarrollo social.
  \item Retroalimentación y nivelación.
\end{itemize}

\section{Lectura editorial}

La bienvenida cumple una función de encuadre inicial: abre la asignatura, presenta al docente, invita a la interacción en sesión sincrónica y establece una pauta básica de convivencia académica. La introducción y el objetivo amplían ese encuadre: la materia no solo crea un ambiente respetuoso, sino que orienta el aprendizaje hacia la defensa profesional de derechos humanos mediante garantías y medios de control constitucional. Para AulaTeX, este contenido funciona como nodo local de memoria de la materia, porque fija tono, actores, conceptos iniciales, ruta de análisis y reglas de respeto que deberán preservarse en futuras actividades.

\section{Conclusión}

El inicio de \textit{Garantías Constitucionales} se organiza alrededor de una comunicación cercana, respetuosa e institucional, pero también alrededor de una finalidad jurídica concreta: comprender derechos humanos, garantías de protección y medios de control constitucional para abordar problemáticas de defensa de derechos. La presencia del docente como abogado con formación de posgrado en la Facultad de Derecho de la UNAM aporta un punto de partida profesional, mientras que el diagnóstico de conocimientos previos permite convertir la bienvenida en una ruta de nivelación, retroalimentación y aprendizaje jurídico aplicable.

\clearpage
\nocite{unadmSitioWeb,unadmMallaDerecho2024,cpeum2026,leyAmparo2026}
\bibliography{garantias-constitucionales}

\end{document}